\section{Secondary diagnostics and additional metrics}

\subsection{H2b secondary analysis}

H2b uses the A--C movement parameterization and is secondary, non-independent,
and unable to rescue H2a. The locally retained Retina outputs give the following
mean changes over 66 class-4 OOF examples.

\begin{center}
\begin{tabular}{llrr}
\toprule
Objective & Seed & Mean $\Delta\mu_p$ & Mean $\Delta p_4$ \\
\midrule
CE & 1 & .225 & .131 \\
CE & 2 & .766 & .254 \\
CE & 3 & .437 & .190 \\
CE & 4 & .821 & .320 \\
RPS & 1 & .388 & .270 \\
RPS & 2 & .648 & .282 \\
RPS & 3 & .635 & .200 \\
RPS & 4 & .660 & .282 \\
\bottomrule
\end{tabular}
\captionof{table}{Retina H2b secondary output.}
\end{center}

\subsection{Global and uncertainty-quality context}

Rare-end localization improvement is not asserted to imply universal global or
uncertainty-quality improvement. The retained historical seed-0 direction
summaries illustrate this separation: global exact-L1 MAE changes from .696 to
.721 for Retina RPS, .693 to .709 for Retina CE, .452 to .552 for Solar RPS,
and .456 to .595 for Solar CE under A to C. The Phase 3.19 summary additionally
retains NLL, Brier, RPS, ECE, risk--error Spearman, severe AUROC/AUPRC, and
selective MAE for A/C/E/F/G. These diagnostics are historical, single-backbone
context rather than an additional confirmatory endpoint.

No new UQ metric, seed aggregation, or selective analysis is introduced here.


\subsection{Natural-sampling direction-only follow-up}

The natural-sampling direction-only analysis is auxiliary and was not part of the frozen H1/H2a confirmatory test. It holds the representation, original class-row norms, and original biases fixed, and changes the C protocol only from replacement-balanced to empirical natural sampling. N is weaker than C in every matched unit.

\begin{center}
\begin{tabular}{lccc}
\toprule
Setting & N improves A & N weaker than C & Descriptive interpretation \\
\midrule
RetinaMNIST CE & 2/4 & 4/4 & Not consistent \\
RetinaMNIST RPS & 2/4 & 4/4 & Not consistent \\
Solar CE & 2/4 & 4/4 & Not consistent \\
Solar RPS & 3/4 & 4/4 & Mixed/partial \\
\bottomrule
\end{tabular}
\captionof{table}{Cross-setting natural-sampling direction-only follow-up. C is the balanced direction-constrained probe.}
\label{tab:natural-followup}
\end{center}

For Solar CE, the four-unit auxiliary summary uses original seeds 1, 3, and 4 plus a pre-declared coherent replacement seed 2. The historical seed-2 features did not meet the frozen numerical identity requirement for this auxiliary follow-up; the original confirmatory record remains preserved. Solar CE mean endpoint MAE is \(1.430/1.426/.538\) for A/N/C, mean global MAE is \(.492/.436/.587\), mean macro MAE is \(.654/.605/.534\), and mean severe-error rate is \(.046/.037/.084\). N shifts endpoint mass only slightly in the upper classes (mean N--A changes \(+.022\) and \(+.038\) for true classes 3 and 4), whereas C produces broader, larger changes (mean C--A changes \(+.201,+.422,+.495\) for true classes 2, 3, and 4).
